% exp-fips203-mlkem-pke-encrypt-k2 — FIPS 203 Alg.14 ML-KEM-512 PKE Encrypt
% 编译：bash scripts/xelatex-clean.sh examples/incubating/ml-kem/ml-kem-512/exp-fips203-mlkem-pke-encrypt-k2/exp-fips203-mlkem-pke-encrypt-k2-实现方案-customspec.tex
\documentclass[UTF8,a4paper,11pt]{ctexart}
\usepackage{amsmath,amssymb}
\usepackage{booktabs}
\usepackage{geometry}
\usepackage{hyperref}
\usepackage{longtable}
\usepackage{array}
\usepackage{seqsplit}
\geometry{margin=2.2cm}
\newcolumntype{L}[1]{>{\raggedright\arraybackslash}p{#1}}
\newcommand{\apiname}[1]{\texttt{\seqsplit{#1}}}
\newcommand{\dirname}[1]{\texttt{\seqsplit{#1}}}
\hypersetup{colorlinks=true,linkcolor=blue,urlcolor=blue}

\title{exp-fips203-mlkem-pke-encrypt-k2\\FIPS 203 Alg.14 ML-KEM-512 PKE Encrypt 实现方案}
\author{ascendc 工程（ML-KEM-512 W4a / E14）}
\date{2026-07-27}

\begin{document}
\maketitle

\begin{abstract}
本文档描述 \dirname{examples/incubating/ml-kem/ml-kem-512/exp-fips203-mlkem-pke-encrypt-k2/}：
在 \textbf{Atlas A2 / Ascend910B4} 上以 AscendC 实现 FIPS 203 \textbf{Algorithm 14}，即 ML-KEM-512（$k=2$）\textbf{PKE Encrypt}。

\textbf{参数锁定}：\dirname{docs/specs/fips203-mlkem512-parameter-card.md} 与
\dirname{ascendc-tests/ml-kem/ml-kem-512/pass-fix-f203-alg14-pke-encrypt-device-k2/STATUS.md}。
\textbf{生产 I/O}：\texttt{input/} 为 \texttt{ek\_pke.bin}（800B）、\texttt{m.bin}（32B）、\texttt{coins.bin}（32B）与静态 LUT；
\texttt{output/} 仅 \texttt{c.bin}（768B）。
\textbf{计划 AI Core 数}：2 launch；Launch 1 为 AIV\_ONLY \texttt{blockDim=2}，Launch 2 为 MIX \texttt{blockDim=1}（S-1）。
\textbf{禁止}：写入 stable-512；从 \dirname{frozen/} 抄码；用 k3/k4 的 \texttt{ek/c} 长度、\texttt{re[7/9]} 或零垫替代 k2 几何。
\end{abstract}

\tableofcontents
\newpage

\section{工程边界}

\begin{longtable}{@{}L{4.6cm}L{8.6cm}@{}}
\toprule
来源 & 用途 \\
\midrule
\dirname{ascendc-tests/ml-kem/ml-kem-512/pass-fix-f203-alg14-pke-encrypt-device-k2/} & 活跃 D14 k2 行为基线；可复制后自包含适配 \\
\dirname{examples/incubating/ml-kem/ml-kem-768/exp-fips203-mlkem-pke-encrypt-k3/} & 活跃 k3 incubating 结构参考；仅 retarget 到 k2 \\
\dirname{library/shared/} & SHAKE/Keccak、统一整数压缩辅助等 \\
CANN / tikicpulib / CAModel & 编译与运行环境 \\
\bottomrule
\end{longtable}

\subsection{禁止项}
\begin{itemize}
  \item 禁止从 \dirname{frozen/} 拿源码、customspec 或路线。
  \item 禁止将 \(\hat A\)、\texttt{re}、\(u\)、\(v\) 等中间态作为默认 I/O 落盘。
  \item 禁止把噪声或 INTT 输入零垫到 6/8；\texttt{re} 锁定为 5 poly。
  \item 禁止新建或写入 stable-512。
\end{itemize}

\section{数学语义（Alg.14，$k=2$，$n=256$，$q=3329$）}

\subsection{输入输出}

\begin{longtable}{@{}L{3.3cm}L{10.2cm}@{}}
\toprule
张量 / 文件 & 契约 \\
\midrule
\texttt{input/ek\_pke.bin} & 800B；前 768B 为 \(\hat t\) 的 \(\mathrm{ByteEncode}_{12}\)，末 32B 为 \(\rho\) \\
\texttt{input/m.bin} & 32B 明文消息 \\
\texttt{input/coins.bin} & 32B Alg.14 随机性；用于 PRF 生成 \(r\Vert e_1\Vert e_2\) \\
\texttt{input/lut\_even\_stacked.bin}、\texttt{lut\_odd\_stacked.bin} & NTT/INTT Stage2 LUT \\
\texttt{output/c.bin} & 768B，\(c=c_1\Vert c_2\)，其中 \(c_1=640\)B（\(d_u=10\)，2 poly），\(c_2=128\)B（\(d_v=4\)，1 poly） \\
\bottomrule
\end{longtable}

\subsection{阶段公式}

\begin{enumerate}
  \item 行 2：从 \texttt{ek\_pke[0:768]} 解码 \(\hat t\in R_q^2\)，在 compute launch 内完成。
  \item 行 3--7：取 \(\rho=\texttt{ek\_pke[768:800]}\)，对 \((j,p)\in\{0,1\}^2\) 生成 \(\hat A[j,p]\)，GM 锁定为 \texttt{[4,256] int32}。
  \item 行 8--15：用 \(\eta=2\) 从 \texttt{coins} 采样 \(r\Vert e_1\Vert e_2\) 共 5 个 poly；GM 锁定为 \texttt{[5,256]}。
  \item 行 16--18：\(\hat r=\mathrm{NTT}(r)\)；\(\hat u_j=\sum_p \hat A[j,p]\circ\hat r_p\)；\(\widehat{tr}=\sum_p \hat t_p\circ\hat r_p\)。
  \item 行 19--21：对 \(\hat u_0,\hat u_1,\widehat{tr}\) 与一空槽执行真 \textbf{INTT batch=4}；得到 \(u_0,u_1,v\) 后分别加 \(e_1\) 与 \(e_2+\mu(m)\)。
  \item 行 22：\(c_1=\mathrm{ByteEncode}_{10}(\mathrm{Compress}_{10}(u))\)，\(c_2=\mathrm{ByteEncode}_{4}(\mathrm{Compress}_{4}(v))\)，输出 \(c=c_1\Vert c_2\)。
\end{enumerate}

\section{AscendC 实现规划}

\subsection{Launch 与 AI Core 规模}

\begin{longtable}{@{}L{2.5cm}L{3.0cm}L{7.6cm}@{}}
\toprule
Launch & 核类型 / blockDim & 数据划分与理由 \\
\midrule
1 prep & AIV\_ONLY，\texttt{blockDim=2} & \(\hat A\) 共 4 poly 按 2+2；$r\Vert e_1\Vert e_2$ 共 5 poly 写真实 GM。继承 D14 k2 已验。 \\
2 compute+tail & MIX，\texttt{blockDim=1} & S-1：NTT(\(r\)) k=2 → 内积 → INTT batch4（槽 \(u_0,u_1,v,\)空）→ Compress/ByteEncode \(d=10/4\)。 \\
\bottomrule
\end{longtable}

\subsection{GM 几何}

\begin{longtable}{@{}L{3.7cm}L{9.5cm}@{}}
\toprule
对象 & 锁定几何 \\
\midrule
\(\hat A\) & \texttt{[4,256] int32}；禁止 pad 到 9/16 \\
\texttt{re} & \texttt{[5,256] int32}；\texttt{r[0..1] || e1[0..1] || e2[0]} \\
\(\hat r\) NTT & polyvec2；AIV 1+1 \\
INTT batch4 & S0 \texttt{[8,256] int8}，mat\_c \texttt{[32,128] int32}；语义槽 \(u0,u1,v,\)空 \\
Tail pack & \(u\) 两 poly \(d_u=10\) → 640B；\(v\) 一 poly \(d_v=4\) → 128B \\
\bottomrule
\end{longtable}

\section{AscendC API 表}

本实现不引入 E14 独有的新 AscendC API；约束按 API 查阅索引执行。

\begin{longtable}{@{}L{3.2cm}L{2.7cm}L{2.5cm}L{2.2cm}L{3.0cm}@{}}
\toprule
API 名 & 分类 / 文档 & Atlas A2 & 本用例 dtype & 备注 \\
\midrule
\apiname{GetBlockIdx} / \apiname{GetSubBlockIdx} & 系统变量 / 资源 & √ & 标量索引 & 区分 AIV 分片与 MIX 子核 \\
\apiname{GlobalTensor} / \apiname{LocalTensor} & 基础结构 & √ & \texttt{uint8/int8/int32} & GM/UB 张量契约须与本 spec 几何一致 \\
\apiname{DataCopy} & Memory 数据搬运 & √ & \texttt{uint8/int8/int32} & 连续块搬运；32B 对齐约束 \\
\apiname{Duplicate} & Memory 矢量计算 & √ & \texttt{uint8/int32} & UB 清零、pad 与 tail buffer 初始化 \\
\apiname{PipeBarrier} & Pipe 同步 & √ & — & MTE/Vector/Cube 阶段交界显式同步 \\
\apiname{CrossCoreSetFlag} / \apiname{CrossCoreWaitFlag} & MIX 同步 & √ & — & AIV/AIC 通过 GM 交接；CPU 过不能删 \\
\apiname{Mmad} / \apiname{Fixpipe} & 矩阵计算 & √ & \texttt{int8*int8->int32} & Stage2 MMAD；硬件 M 对齐垫不表示假 poly \\
\apiname{ShiftRight} & 矢量算术 & √ & \texttt{int32} & limb 拆分、Barrett 与 ByteEncode 辅助右移 \\
\apiname{Muls} / \apiname{Add} / \apiname{Sub} & 矢量算术 & √ & \texttt{int32/int16} & limb、模加减、压缩/编码辅助 \\
\apiname{Cast} & 类型转换 & √（受限） & \texttt{int32/int16/half/int8} & 遵守 A2 Cast 链；不得假设 \texttt{int32->int8} 单跳 \\
\apiname{Gather} & 矢量搬运/选择 & √ & \texttt{int32} & 仅允许非 NTT S1--S3 段；NTT 三段内禁用 \\
\apiname{GetValue} / \apiname{SetValue} & LocalTensor 标量访问 & √ & \texttt{uint8/int32} & 仅用于 pack 尾部等标量缺口 \\
\bottomrule
\end{longtable}

\subsection{评估过但未采用}

\begin{longtable}{@{}L{3.8cm}L{9.2cm}@{}}
\toprule
API / 路线 & 不采用原因 \\
\midrule
零垫 \texttt{re} 到 7/9 & T-B2 / 禁零垫；D14 k2 已锁 5 poly \\
\apiname{Gather} 用于 NTT S1--S3 & NTT 定稿禁止 \\
高阶 \apiname{Matmul<>} & 继承 D14 手写 MMAD + 平面 mat\_c \\
\bottomrule
\end{longtable}

\section{数学步骤覆盖率}

\begin{longtable}{@{}L{4.2cm}L{4.0cm}L{2.0cm}L{3.0cm}@{}}
\toprule
数学步骤 & 实现方式 / API & 覆盖 & 缺口说明 \\
\midrule
解码 \(\hat t\)、取 \(\rho\) & compute 内 ByteDecode12 + 指针切片 & 部分向量 & unpack 标量缺口沿用 D14 \\
采样 \(\hat A[4]\) & AIV\_ONLY + XOF/rej & 100\% 设备 & 2+2；无零垫 \\
采样 \(r\Vert e_1\Vert e_2[5]\) & AIV\_ONLY + CBD \(\eta=2\) & 100\% 设备 & 真实 5 poly \\
NTT / 内积 / INTT batch4 & MIX + \apiname{Mmad} & 100\% 设备 & NTT 内无 Gather；空槽不写假 poly 语义 \\
Compress/ByteEncode \(d=10/4\) & 统一整数压缩 + pack & 部分向量 & pack 尾部标量缺口沿用 D14 \\
\bottomrule
\end{longtable}

\section{验证计划}

\begin{verbatim}
cd examples/incubating/ml-kem/ml-kem-512/exp-fips203-mlkem-pke-encrypt-k2
bash run.sh -r cpu -v Ascend910B4
SIM_DIRECT=1 bash run.sh -r sim -v Ascend910B4
\end{verbatim}

期望：\texttt{c.bin} 为 768B，与 golden 对拍 PASS；用例根无 stray dump。

\section{产出物}
\begin{itemize}
  \item \dirname{exp-fips203-mlkem-pke-encrypt-k2-实现方案-customspec.tex}
  \item \dirname{exp-fips203-mlkem-pke-encrypt-k2-实现方案-customspec.pdf}
  \item 后续实现文件须与本 customspec 同目录自包含，且自研 Python/C/AscendC 写详细中文注释。
\end{itemize}

\end{document}
